\documentclass[17pt,a4paper,landscape]{extarticle}
\usepackage[margin=2.35cm,top=0.12cm,left=1.05cm]{geometry}
\usepackage{amsmath}
\usepackage{amssymb}
\usepackage{tasks}
\usepackage{tikz}
\usepackage{xcolor}
\usepackage[sfdefault,lf]{FiraSans}
\renewcommand{\familydefault}{\sfdefault}
\setlength{\parindent}{0pt}
\pagestyle{empty}
\settasks{label=\textbf{\arabic*.}, label-width=2.2em, item-indent=2.95em, column-sep=2.1cm, after-item-skip=4.7em}
\renewcommand{\arraystretch}{1.15}
\everymath{\displaystyle}

\begin{document}
\sffamily
\boldmath

\boldmath

\boldmath

\boldmath

{\LARGE \textbf{Sampling methodology --- Excellence}}
\begin{tasks}(3)
\task A principal wants to estimate support for a new uniform. Compare surveying students outside the uniform shop with a stratified random sample across year levels.
\task Explain the difference between a large sample and a representative sample.
\task A population is \mbox{$55\%$} juniors and \mbox{$45\%$} seniors. Design a stratified sample of size \mbox{$60$}.
\task A screen-time survey is sent through the school e-sports club only. Give one criticism using selection bias.
\task A researcher takes every \mbox{$20$}th name from an alphabetical roll. Give one strength and one weakness of this method.
\task Two surveys both use \mbox{$100$} people. One is random across the town. The other uses \mbox{$100$} shoppers from one mall at midday. Which is stronger? Why?
\task A sports survey includes \mbox{$30$} basketball players, \mbox{$25$} netball players, and only \mbox{$2$} students who do no sport. Why might this be unrepresentative?
\task A researcher says, ``My sample is big, so it must be fair.'' Explain why that does not always follow.
\task A company wants opinions on a new app from teenagers and adults. Why might stratifying by age help?
\task One sample is random but too small. Another is larger but clearly biased. Which problem is harder to fix later? Explain.
\task A town has \mbox{$40\%$} urban, \mbox{$35\%$} suburban, and \mbox{$25\%$} rural households. For a stratified sample of \mbox{$200$}, how many come from each group?
\task A researcher chooses every \mbox{$10$}th customer from a list arranged by membership type. Why could this be biased?
\task Why do volunteer surveys often over-represent people with strong opinions?
\task Recommend a good sampling methodology for estimating travel-to-school patterns across a whole school.
\task A survey uses a perfect random method, but half the selected people refuse to answer. What new concern does this create?
\task Give an example where a cluster-like convenience sample could be useful, but risky for generalising.
\end{tasks}

\end{document}
